%% ma: L12-B08-0075
%% lop: 12
%% chuong: 2
%% bai: 8
%% dang: 12.8.30
%% loai: TLN
%% muc: VD
%% dap_an: "1"
%% nguon: "LaTeX giáo viên gửi 21/10/2026 (M05, Vectơ và hệ trục tọa độ trong không gian, 107 câu), Phần III Câu 15"
%% kien_thuc: "Bài 6 (quy tắc trung điểm), Bài 7 (hình chiếu lên mặt phẳng tọa độ), Bài 8 (trung điểm)"
%% trang_thai: cho-duyet
%% ly_do: "Dạng mới đề xuất 12.8.30"
%% ghi_chu: "Bỏ tiêu đề '(Bài toán Cực trị)'"
\begin{ex}
Trong không gian $Oxyz$, cho hai điểm $A(1;-2;4)$ và $B(3;0;-2)$. Coi mặt phẳng $(Oxy)$ là mặt đất. Cần chọn một vị trí $M(a;b;0)$ trên mặt đất sao cho $|\overrightarrow{MA}+\overrightarrow{MB}|$ đạt giá trị nhỏ nhất. Tính $S=a+b$.
\shortans{1}
\loigiai{Gọi $I(2;-1;1)$ là trung điểm $AB$, $\overrightarrow{MA}+\overrightarrow{MB}=2\overrightarrow{MI}$. $|\overrightarrow{MA}+\overrightarrow{MB}|$ nhỏ nhất khi $M$ là hình chiếu của $I$ lên $(Oxy)$: $M(2;-1;0)$, $S=1$.}
\end{ex}
